\documentclass{article}
% empty [] options -> geometry falls back to its default margins here
\usepackage[]{geometry}
\usepackage{graphicx}
\usepackage{float}
\usepackage{wrapfig}
\usepackage{caption}
\usepackage{subcaption}
\usepackage{booktabs}
\usepackage{array}
\usepackage{enumitem}
\usepackage[table]{xcolor}
\usepackage{amsmath, amssymb, amsfonts}
\usepackage{listings}
\usepackage{algorithm}
\usepackage{algpseudocode}
\usepackage{comment}
\usepackage[numbers]{natbib}
\usepackage{hyperref}
\usepackage{subcaption}
\usepackage{multirow}
\usepackage{amsthm}
\usepackage{minitoc}
\usepackage{authblk}

\title{\textsc{CSE200: Online-1 on \LaTeX}}
\author{Zarif Ishmam (2305035)}
\date{\today}

\begin{document}

\maketitle

\section{Introduction}
Document preparation systems are essential for producing structured and professional technical documents. \textbf{LATEX} allows authors to focus on content
while maintaining consistent formatting. This article demonstrates formatted
text, lists, equations, tables, and figures using commands taught in class.
\subsection{Text Emphasis and Fonts}
This sentence contains \textbf{bold text}, \textit{italic text}, \emph{emphasized text}, and \underline{underlined text}.
Font sizes can also vary such as {\LargeLarge text}, {\smallsmall text}, and 
{\huge huge text}.
\section{ List Structures}
\subsection{Mixed and Nested Lists}
% itemize -> enumerate -> itemize, three levels deep
\begin{itemize}
    \item Primary Feature
    \item Secondary Feature
        \begin{enumerate}
            \item Formatting Control
            \item Mathematical typesetting
                \begin{itemize}
                    \item[-] Inline Math
                    \item[-] Display Math
                \end{itemize}
        \end{enumerate}
\end{itemize}
\noindent \textbf{Note} \text{Lists can be customized.}

% note: this "Mathematical Modeling" section is duplicated in the source article itself (title repeats below), kept as-is to match it exactly
\section{Mathematical Modeling}
Mathematical expressions are often used to represent data relationships. Consider variables $x_i$, $y$, and $z^2$.

\section{Mathematical Modeling}
Mathematical expressions are commonly used to describe relationships among variables. Let the variables be $x_i$, $y_j$, and $\sigma^2$.

\subsection{Displayed Equations}
% single-line numbered equation
\begin{equation}
    y_j = \alpha x^2 + \beta x_j + \gamma
\end{equation}

% split keeps this multi-line derivation under one shared equation number
\begin{equation}
\begin{split}
    T &= \sum_{i = 1}^{n}\left( x_i - \mu \right) ^2 \\
    &= \sum_{i = 1}^{n}x_i^2 - 2\mu \sum_{i = 1}^{n} x_i + n \mu ^2
\end{split}
\end{equation}

% cases for the piecewise definition (unnumbered, via \[ \])
\[
f(x_i) = \begin{cases}
    \frac{x_i^2}{\sigma^2},  \text{if} x_i >= 0 \\
    \frac{|x_i|}{\sigma}, \text{if} x_i < 0
\end{cases}
\]

\section{Tabular Data Presentation}
% multirow for "Module"/"Module B", multicolumn + cline for the grouped "Score" header
\begin{table}[H]
    \centering
    \begin{tabular}{|c|c|c|}
        \hline
         \multirow{2}{*}{\textbf{Module}} & \multicolumn{2}{c|}{\textbf{Score}}  \\ \cline{2-3} & Theory & Lab \\
         \hline
         Module A & 75 & 80 \\
         \hline
         \multirow{2}{*}{Module B} & 88 & 90 \\ & 85 & 87 \\
         \hline
         Module C & 70 & 72 \\
         \hline
    \end{tabular}
    \caption{Module-wise Performance Summary}
    \label{tab:t1}
\end{table}

\section{Visual Comparison}
\subsection{Asymmetric Subfigure Layout}
% one wide subfigure on its own row, then two narrower ones side by side below

\begin{figure}[H]
    \centering
    \begin{subfigure}{0.6\textwidth}
        \centering
        \includegraphics[width=\linewidth]{22_C2/A.PNG}
        \caption{Main Diagram}
        \label{fig:A}
    \end{subfigure}
    \par\medskip
    \vspace{1cm}
    \begin{subfigure}{0.3\textwidth}
        \centering
        \includegraphics[width=\linewidth]{22_C2/B.PNG}
        \caption{Main Diagram}
        \label{fig:B}
    \end{subfigure}
    \hfill
    \begin{subfigure}{0.3\textwidth}
        \centering
        \includegraphics[width=\linewidth]{22_C2/C.PNG}
        \caption{Main Diagram}
        \label{fig:C}
    \end{subfigure}
    \caption{Asymmetric layout with one dominant figure}
    \label{fig:f1}
\end{figure}

Figure \ref{fig:f1} illustrates how visual elements can be arranged systematically within a document.

\section*{Conclusion}
This article highlighted the importance of \textbf{structured writing} using \LaTeX.
The use of \textit{formatted text}, mathematical expressions, tables, and figures improves both \underline{readability} and presentation quality.

\end{document}